\subsection{Step-by-Step Procedure}\label{sec:procedure}
Algorithm~\ref{alg:dc} restates Algorithm~2.8 of~\cite{lunagarcia2024performance}.

\begin{algorithm}[H]
\caption{Divide and Conquer approach for TSP (Algorithm~2.8 of~\cite{lunagarcia2024performance})}\label{alg:dc}
\begin{algorithmic}[1]
\Require A set of cities with coordinates and pairwise distances
\State Randomly shuffle the cities to introduce variation.
\State Recursively split the cities into two subsets by their $x$ or $y$ coordinate (chosen randomly).
\If{the subset has 1 city} \State return the city
\ElsIf{the subset has 2 cities} \State return the edge connecting the two cities
\EndIf
\State Recursively solve the TSP for each subset.
\State Merge the two solutions through the shortest connection (bridge) between the subsets.
\State Adjust the merged solution so that no city is repeated.
\Ensure The merged tour and its total distance
\end{algorithmic}
\end{algorithm}

\paragraph{Step 1: Shuffle.} The input order is randomised. Our split sorts by coordinate using a
deterministic coordinate-and-index key, so the shuffle does not change the splits. It is kept to follow
the paper, and it only advances the random-number generator; the randomness that changes the result is
the choice of axis in Step~2.

\paragraph{Step 2: Divide.} Pick the $x$ axis (longitude) or the $y$ axis (latitude) at random, sort the
cities along it, and cut at the median. This gives two halves of sizes $\lfloor m/2\rfloor$ and
$\lceil m/2\rceil$ whose points lie on opposite sides of a vertical or horizontal line.

\paragraph{Step 3: Base cases.} A single city is a ``tour'' of length 0. Two cities $a,b$ form the tour
$a \to b \to a$ of length $2w(a,b)$; the paper's code uses the same value.

\paragraph{Step 4: Conquer.} Solve each half recursively, which returns a cycle $T_1$ for the first half
and $T_2$ for the second.

\paragraph{Steps 5--6: Combine (merge).} The two cycles must become \emph{one} cycle that still visits
each city once. The paper's code (Appendix~4.8 of~\cite{lunagarcia2024performance}) concatenates
\texttt{tour1[:-1] + [a, b] + tour2}, where $(a,b)$ is the shortest bridge. This does not give a valid
tour, because $a$ and $b$ are already in \texttt{tour1} and \texttt{tour2}. A preserved run of this
paper-style expression on a seven-city subproblem is stored in
\texttt{code/original\_merge\_output.txt}: the returned sequence has eight entries, repeats V5, and
does not close at its start. In other words, the step ``adjust the merged solution so no city is repeated''
is never actually carried out.

We therefore implement the merge in the standard way for joining two cycles, which keeps the paper's
``shortest bridge'' rule:

\begin{algorithm}[H]
\caption{Merge two cycles through the shortest bridge (our implementation of Steps 5--6)}\label{alg:merge}
\begin{algorithmic}[1]
\Require Cycles $T_1$ and $T_2$ on disjoint vertex sets $S_1$ and $S_2$; weights $w$
\State $(a, b) \gets \arg\min_{u \in S_1,\, v \in S_2} w(u,v)$ \Comment{shortest bridge}
\For{each edge $(a, a')$ of $T_1$ at $a$ and each edge $(b, b')$ of $T_2$ at $b$}
  \State remove $(a,a')$ and $(b,b')$; add the bridge $(a,b)$ and the closing edge $(a',b')$
  \State $\Delta \gets w(a,b) + w(a',b') - w(a,a') - w(b,b')$
\EndFor
\State keep the option with the smallest $\Delta$; the merged cycle has length $L(T_1)+L(T_2)+\Delta$
\Ensure One cycle through $S_1 \cup S_2$ that visits each vertex exactly once
\end{algorithmic}
\end{algorithm}

Removing one edge from each cycle leaves two paths; the bridge and the closing edge join their ends into
one cycle, so every city appears exactly once. A 1-city tour has no edge to remove, so the city is simply
inserted next to $b$. For a 2-city tour $a \to b \to a$, removing one copy of the doubled edge leaves the
single edge $a$--$b$. Section~\ref{sec:apply} shows this procedure on our data.

\begin{remark}
We also implemented a variant, the \emph{best-splice} rule, that does not insist on using the shortest
bridge: it tries every edge of $T_1$ against every edge of $T_2$ and keeps the smallest $\Delta$. It is
still $O(n^2)$ per merge. Unless stated otherwise, all results use the paper-faithful shortest-bridge
rule; the variant is compared in Sections~\ref{sec:optimality} and~\ref{sec:limitations}.
\end{remark}
